\documentclass[11pt,a4paper]{article}

\usepackage[margin=2.5cm]{geometry}
\usepackage[T1]{fontenc}
\usepackage[utf8]{inputenc}
\usepackage{lmodern}
\usepackage{microtype}
\usepackage{booktabs}
\usepackage{amsmath}
\usepackage{xcolor}
\usepackage{hyperref}
\usepackage{listings}

\hypersetup{colorlinks=true, linkcolor=blue!50!black, urlcolor=blue!60!black}
\lstset{
  basicstyle=\ttfamily\small,
  frame=single,
  breaklines=true,
  columns=fullflexible,
  backgroundcolor=\color{black!3}
}

\title{Personalized Hand-Gesture Recognition\\
\large Classical ML Models Compared with Hand-Written Rules}
\author{Gesture Presenter --- Machine Learning Course Study}
\date{13 July 2026}

\begin{document}
\maketitle

\begin{abstract}
Gesture Presenter recognizes hand gestures using MediaPipe landmarks and
hand-written geometric rules. This study tested whether classical machine
learning could classify one user's gestures more accurately. A decision tree,
random forest, and logistic regression model were trained on two recording
sessions and tested on a separate third session. Only hand landmarks were
stored; no images or video were saved. The random forest performed best with
98.5\% trial accuracy, followed by logistic regression with 94.8\%, a decision
tree with 93.3\%, and the existing rules with 79.3\%. The main improvement was
recognizing natural, non-command movement instead of incorrectly treating it as
a command. This is a personalized single-user result and does not show that the
models generalize to other people.
\end{abstract}

\section{Goal}

Gesture Presenter uses MediaPipe to detect 21 landmarks on a hand. The current
application then recognizes gestures with manually selected thresholds, such as
finger joint angles and distances between fingertips. These rules are fast and
easy to understand, but the same fixed thresholds may not suit every user's hand
shape and gesture style.

The research question was:

\begin{quote}
Can a small, explainable classical ML model recognize one user's static gestures
more accurately than the existing hand-written rules across separate recording
sessions?
\end{quote}

The aim was not to replace MediaPipe. MediaPipe remained responsible for hand
detection. The experiment compared two ways of mapping the same landmarks to a
gesture:

\begin{center}
Camera $\rightarrow$ MediaPipe landmarks $\rightarrow$
\{hand-written rules or trained model\} $\rightarrow$ gesture
\end{center}

This was designed as a personalized pilot because only one participant and one
evening were available. Therefore, the study measures repeatability for one
user, not generalization to the public.

\section{Data collection}

A guided collection script displayed a gesture, a countdown, a 1.5-second
recording period, and a rest period. Gesture order was randomized. The
participant used the right hand and recorded three sessions:

\begin{table}[ht]
\centering
\caption{Dataset and experimental split.}
\begin{tabular}{lrrl}
\toprule
Session & Frames & Grouped trials & Use \\
\midrule
S01 & 5,284 & 135 & Training \\
S02 & 3,509 & 135 & Training \\
S03 & 4,762 & 135 & Final evaluation \\
\midrule
Total & 13,555 & 405 & --- \\
\bottomrule
\end{tabular}
\end{table}

Nine classes were recorded:

\begin{center}
\begin{tabular}{lll}
\toprule
\texttt{none} & \texttt{point} & \texttt{pinch} \\
\texttt{pinch\_right} & \texttt{open\_palm} & \texttt{fist} \\
\texttt{fist\_thumb\_up} & \texttt{v\_sign} & \texttt{two\_finger\_scroll} \\
\bottomrule
\end{tabular}
\end{center}

The \texttt{none} class contained natural hand movement without an intended
command. It is important because false commands are particularly disruptive in
a presentation controller.

Camera frames were processed locally and immediately discarded. Each saved CSV
row contained the participant, session and trial identifiers, gesture label,
timestamp, handedness information, and the $(x,y,z)$ coordinates of all 21
landmarks. No original footage was stored. Landmark data is less sensitive than
video, although it should not be described as completely anonymous.

\section{Features and models}

Each frame was represented by 77 features:

\begin{itemize}
  \item 63 coordinates: 21 landmarks with $x$, $y$, and $z$ values;
  \item four normalized distances between important hand points;
  \item horizontal and vertical thumb displacement;
  \item four finger joint angles; and
  \item four fingertip-to-wrist distances.
\end{itemize}

Coordinates were translated so the wrist became the origin and divided by hand
size. This reduces sensitivity to where the hand appears in the image and how
close it is to the camera.

Three models were trained with balanced class weights and random seed 42:

\begin{description}
  \item[Decision tree:] maximum depth 8 and at least 5 frames per leaf. It was
  selected first because each prediction can be explained as a short sequence
  of feature thresholds and compared directly with the manual rules.

  \item[Random forest:] 250 trees and at least 3 frames per leaf. Combining many
  trees normally improves stability, but makes individual decisions less direct
  to explain.

  \item[Logistic regression:] standardized features and up to 2,000 training
  iterations. This tests whether simpler, approximately linear boundaries are
  sufficient for the engineered features.
\end{description}

The existing geometric classifier was the baseline. It uses finger-extension
tests, joint angles, fingertip distances, thumb direction, and manually chosen
thresholds.

\section{Evaluation}

All learned models were trained on 8,793 frames from S01 and S02. S03 was
excluded from training and used only for evaluation. The comparison program
checks model metadata and refuses to evaluate a session that was used for
training. All classifiers received the same S03 landmarks in the same order.

Two accuracy measures were calculated:

\begin{description}
  \item[Frame accuracy] scores every camera frame independently.
  \item[Trial accuracy] combines all frame predictions in one prompted
  repetition using majority voting.
\end{description}

Trial accuracy is the main measure because adjacent frames are highly similar
and are not independent samples. One complete gesture repetition is a more
meaningful unit than one frame.

For each gesture, precision and recall describe different types of error:

\[
\text{Precision} = \frac{TP}{TP+FP},
\qquad
\text{Recall} = \frac{TP}{TP+FN},
\]

where $TP$ is the number of correct predictions of that gesture, $FP$ is the
number of other gestures incorrectly predicted as that gesture, and $FN$ is the
number of real examples of that gesture that the classifier missed. Precision
therefore asks, ``When the model predicts this gesture, how often is it
correct?'' Recall asks, ``Of all real examples of this gesture, how many did the
model find?''

The F1 score combines precision and recall using their harmonic mean:

\[
F1 = 2\cdot
\frac{\text{Precision}\cdot\text{Recall}}
{\text{Precision}+\text{Recall}}.
\]

F1 ranges from 0 to 1. A value of 1 means that a gesture was found every time
and was never predicted incorrectly. A low value means that precision, recall,
or both were poor. F1 is useful here because accuracy alone can hide a serious
problem with one gesture class. For example, a classifier may be accurate
overall while frequently missing pinches or treating natural movement as a
command.

Macro F1 is the unweighted mean of the F1 scores for all $K$ gesture classes:

\[
\text{Macro F1} = \frac{1}{K}\sum_{k=1}^{K} F1_k.
\]

Every gesture contributes equally to macro F1, regardless of how many examples
it has. This makes it appropriate for this dataset because the held-out class
supports are not exactly equal.

This comparison covers static classification only. The live application also
uses hold durations, debouncing, swipe tracking, scrolling trajectories, and
cooldowns. Therefore, the experiment does not yet measure actual unintended
clicks or slide changes.

\section{Results}

\subsection{Overall comparison}

\begin{table}[ht]
\centering
\caption{Results on the held-out S03 session.}
\begin{tabular}{lrrr}
\toprule
Classifier & Frame accuracy & Trial accuracy & Correct trials \\
\midrule
Random forest & \textbf{98.8\%} & \textbf{98.5\%} & 133/135 \\
Logistic regression & 95.7\% & 94.8\% & 128/135 \\
Decision tree & 91.7\% & 93.3\% & 126/135 \\
Hand-written rules & 80.0\% & 79.3\% & 107/135 \\
\bottomrule
\end{tabular}
\end{table}

All three learned models outperformed the rules. The random forest improved
trial accuracy by 19.2 percentage points and made only two trial-level errors.
The single decision tree was less accurate than the forest but remains much
easier to explain because it uses one visible path of threshold decisions.

\subsection{Performance by gesture}

Table~\ref{tab:f1} summarizes trial-level F1 scores. F1 combines precision
(how often a prediction is correct) and recall (how many real examples are
found).

\begin{table}[ht]
\centering
\caption{Trial-level F1 score by gesture.}
\label{tab:f1}
\small
\begin{tabular}{lrrrr}
\toprule
Gesture & Forest & Logistic & Tree & Rules \\
\midrule
fist & 1.00 & 1.00 & 0.97 & 0.73 \\
fist\_thumb\_up & 1.00 & 0.89 & 0.96 & 0.89 \\
none & 1.00 & 0.93 & 0.90 & 0.00 \\
open\_palm & 1.00 & 1.00 & 1.00 & 0.84 \\
pinch & 0.97 & 0.88 & 0.77 & 0.83 \\
pinch\_right & 0.97 & 0.93 & 0.90 & 0.85 \\
point & 0.96 & 0.96 & 0.96 & 0.63 \\
two\_finger\_scroll & 0.97 & 0.94 & 0.97 & 1.00 \\
v\_sign & 1.00 & 1.00 & 0.96 & 0.92 \\
\midrule
Macro average & 0.99 & 0.95 & 0.93 & 0.67 \\
\bottomrule
\end{tabular}
\end{table}

The most important result was the \texttt{none} class:

\begin{itemize}
  \item the forest and decision tree recognized all 14 \texttt{none} trials;
  \item logistic regression recognized 13 of 14;
  \item the rules recognized 0 of 14.
\end{itemize}

The rules interpreted the 14 natural-movement trials as seven open palms, three
fists, two pinches, one point, and one thumb-left gesture. These would be
potential false commands before temporal filtering is applied.

The random forest made only two errors: one true index pinch was predicted as
point, and one two-finger-scroll pose was predicted as right-click pinch. The
single decision tree's main weakness was index pinch. It recognized 12 of 19
pinch trials; the other seven were distributed across five other classes. The
forest increased pinch recall from 0.63 to 0.95, showing the benefit of combining
multiple trees.

Logistic regression produced strong results with a simpler decision function.
Its mistakes were mainly pinch and two-finger-scroll examples. The numerical
backend printed overflow warnings during fitting, but the input features,
trained coefficients, predictions, and model output were checked and were all
finite. The training completed in 78 iterations. This score should still be
reproduced in a clean environment before being treated as a final benchmark.

\section{What the results mean}

For this participant, learned decision boundaries transferred from two sessions
to a separate third session better than the fixed rules. The random forest was
the best predictive model. The decision tree remains the best choice for direct
XAI because its learned thresholds can be read and compared with the existing
manual thresholds.

The strongest product-relevant finding is not only higher accuracy. It is the
ability to reject natural hand movement. The current rules are good at several
clearly posed gestures, including two-finger scroll, but often force ordinary
movement into a command class. The learned models captured the participant's
non-command poses much better.

The valid conclusion is:

\begin{quote}
For one participant, three classical ML models trained on two sessions achieved
higher static gesture-classification accuracy than the existing rules on a
held-out third session. The random forest was most accurate, while the single
decision tree offered the clearest explanations.
\end{quote}

It would not be valid to claim that these models work for arbitrary users or
that the random forest is already better in the complete live application.

\section{Limitations and next steps}

\begin{enumerate}
  \item Only one person participated, so there is no cross-user evidence.

  \item Adjacent frames are correlated. Trial-level results are more meaningful
  than the larger frame count.

  \item The class supports in S03 range from 11 to 19 instead of exactly 15.
  This is consistent with restarting randomized collection using the same
  session and trial identifiers, which can mix frames inside a grouped trial.
  A final confirmation session should use a fresh session ID and must not be
  used to tune the existing models.

  \item The models were compared using fixed configurations. S03 must not now be
  reused for hyperparameter selection, because that would leak test information.

  \item The next XAI step is to visualize the decision tree, inspect feature
  importance, and explain misclassified pinch trials. The next product-level
  experiment should apply temporal filtering and measure false actions per
  minute.
\end{enumerate}

\appendix
\section{Reproduction}

The model type can be \texttt{tree}, \texttt{forest}, or \texttt{logistic}:

\begin{lstlisting}
venv/bin/python -m research.train \
  research/data/landmarks.csv \
  --model forest \
  --exclude-session S03 \
  --output research/models/forest.joblib

venv/bin/python -m research.compare \
  --model research/models/forest.joblib \
  --session S03 \
  research/data/landmarks.csv
\end{lstlisting}

\end{document}
